% Standalone orbifold Donaldson--Thomas statements for the fibrewise
% quotients [K3 x E / Z_N] at N=4 and N=6.

\providecommand{\kBKM}{\kappa_{\mathrm{BKM}}}
\providecommand{\kch}{\kappa_{\mathrm{ch}}}
\providecommand{\kcat}{\kappa_{\mathrm{cat}}}
\providecommand{\kfib}{\kappa_{\mathrm{fiber}}}
\providecommand{\DTred}{Z^{\mathrm{red}}_{\mathrm{DT}}}
\providecommand{\FCHL}{F^{\mathrm{CHL}}}
\providecommand{\Z}{\mathbb{Z}}
\providecommand{\mult}{\mathrm{mult}}

\section*{Orbifold Donaldson--Thomas identities at
\texorpdfstring{$N=4$ and $N=6$}{N=4 and N=6}}
\label{wn:sec:N46-orbifold-DT}

Let $S$ be a K3 surface, let $E$ be an elliptic curve, and let
$g_N\in\mathrm{Aut}_s(S)$ be a Mukai--Nikulin symplectic automorphism of
order $N$. The fibrewise orbifold is
\[
  \mathcal{X}^{(N)}=[S\times E/\langle g_N\rangle],
\]
where $g_N$ acts trivially on $E$. Its local stabiliser charts are
$[\mathbb{C}^2/\mathbb{Z}_r]\times E$, equivalently
$[\mathbb{C}^3/\mathbb{Z}_r]$ with the third weight zero. The free
diagonal quotient
\[
  X_N=(S\times E)/\langle(g_N,t_N)\rangle,\qquad
  t_N\in E[N]\ \hbox{of exact order }N,
\]
is a smooth finite etale CHL threefold. The fibrewise orbifold
$\mathcal{X}^{(N)}$ and the smooth CHL quotient $X_N$ carry different
local sectors.

\paragraph{Primitive CHL denominator.}
For the primitive CHL Borcherds ladder
$N\in\{1,2,3,4,6\}$, Borcherds' weight theorem applied to the
$g_N$-twined weight-zero index-one K3 Jacobi input gives
\[
\begin{array}{c|ccccc}
N&1&2&3&4&6\\ \hline
c_N(0)&10&8&6&4&2\\
\kBKM(\Phi_N)&5&4&3&2&1 .
\end{array}
\]
Thus
\[
  \mathrm{wt}\,\FCHL_4=\kBKM(\Phi_4)=2,\qquad
  \mathrm{wt}\,\FCHL_6=\kBKM(\Phi_6)=1 .
\]
The conditional reduced-DT assertions below use the scalar denominator
$(\FCHL_N)^2$, whose weight is twice the primitive Borcherds weight.

The Gritsenko--Cl\'ery diagonal-divisor catalogue is a separate eight-row
dd-modular family with twice-weight tuple
\[
  (10,4,6,2,4,1,3,2).
\]
The level-one row $F_{1,1}=\Delta_5$ is the unique intersection with
the primitive CHL ladder. A comparison with a Gritsenko--Cl\'ery row
requires its own row map, paramodular group, multiplier, branch, and
Borcherds input.

\paragraph{Compact \texorpdfstring{$K3\times E$}{K3 x E} invariants.}
Kunneth gives
\[
  H^{0,\bullet}(S\times E)
  \cong H^{0,\bullet}(S)\otimes H^{0,\bullet}(E),
  \qquad h^{0,q}(S\times E)=(1,1,1,1).
\]
Consequently
\[
  \kcat(S\times E)=\chi(\mathcal{O}_{S\times E})=0,\qquad
  \kch(S\times E)=\sum_q(-1)^q h^{0,q}(S\times E)=0 .
\]
The Heisenberg--Mukai specialisation is a different chiral construction:
\[
  \kch^{\mathrm{Heis}}=3,\qquad
  \kBKM(\mathfrak{g}_{\Delta_5})=\mathrm{wt}(\Delta_5)=5,\qquad
  \kfib=24 .
\]
The four numbers $0,3,5,24$ therefore arise from the compact Hodge
supertrace, the Heisenberg--Mukai specialisation, the Borcherds weight
of $\Delta_5$, and the K3 Mukai-lattice rank.

For $\mathcal{X}^{(N)}$, the same invariant Hodge-column computation
gives
\[
  H^{0,\bullet}(\mathcal{X}^{(N)})
  =
  \bigl(H^{0,\bullet}(S)\otimes H^{0,\bullet}(E)\bigr)^{\langle g_N\rangle}
  =
  (1,1,1,1),
\]
since $g_N$ fixes $H^{0,0}(S)$ and $H^{0,2}(S)$ and acts trivially on
$E$. Hence
\[
  \kcat(\mathcal{X}^{(N)})=0,\qquad
  \kch(\mathcal{X}^{(N)})=0 .
\]

\subsection*{$N=4$, Mathieu class $4B$}

Let $g_4$ have Mathieu cycle shape $1^{4}2^{2}4^{4}$. Nikulin's
fixed-point counts are
\[
  \#\mathrm{Fix}(g_4)=4,\qquad
  \#\mathrm{Fix}(g_4^2)=8 .
\]
The quotient $S/\langle g_4\rangle$ has four stabiliser-$\Z/4$ orbits
and two stabiliser-$\Z/2$ orbits. Its Du-Val type is
\[
  4A_3+2A_1 .
\]
The topological quotient has
\[
  e(S/\Z_4)=\frac{24+4+8+4}{4}=10.
\]
The exceptional curves in the crepant resolution contribute
$4\cdot3+2\cdot1=14$, so the resolved surface has Euler characteristic
$24$.

Let $r_4=\mathrm{rk}\,H^2(S,\mathbb{Z})^{g_4}$. The Lefschetz traces
\[
  \mathrm{tr}(g_4\mid H^2(S))=2,\qquad
  \mathrm{tr}(g_4^2\mid H^2(S))=6
\]
give $r_4=8$ by decomposing $H^2(S,\mathbb{C})$ into
$1,-1,i,-i$ eigenspaces. Therefore the invariant Mukai-lattice rank is
\[
  \mathrm{rk}\,\widetilde H(S,\mathbb{Z})^{g_4}=10 .
\]

\paragraph{Conditional denominator-to-DT assertion at level four.}
Assume a level-four Lorentzian BKM datum with root lattice, Weyl chamber,
real roots, imaginary roots, parity, multiplicities, and denominator
identity is supplied, with primitive automorphic product $\FCHL_4$ and
with root multiplicities
\[
\mult_{\mathfrak{g}^{(4)}}(n\delta+\ell z+mw)
\;=\;c^{(g_4)}_{K3}\!\bigl(4nm-\ell^2,\ell\bigr)
\]
in the positive chamber. Assume also a reduced orbifold-DT gluing theorem
identifying the local $A_3$ and $A_1$ vertex factors, the reduced
$S\times E$ cosection, and the primitive CHL denominator normalisation.
Under these hypotheses,
\begin{equation}
\DTred\bigl(\mathcal{X}^{(4)}\bigr)(\tau,\sigma,z)
\;=\;-\,\frac{C_4(\tau,z)}{\bigl(\FCHL_{4}(\tau,\sigma,z)\bigr)^2},
\qquad
C_4=2\phi^{(g_4)}_{0,1},\quad
\mathrm{wt}\,\FCHL_4=c_4(0)/2=2 .
\label{wn:eq:N4-DTred}
\end{equation}
Here $\phi^{(g_4)}_{0,1}\in J_{0,1}(\Gamma_0(4))$ is the $4B$-twined K3
weak Jacobi form and $c_4(0)=4$.

\subsection*{$N=6$, Mathieu class $6A$, Niemeier anchor $6D_4$}

Let $g_6$ have Mathieu cycle shape $1^{2}2^{2}3^{2}6^{2}$. Nikulin's
fixed-point counts are
\[
  \#\mathrm{Fix}(g_6)=2,\qquad
  \#\mathrm{Fix}(g_6^2)=6,\qquad
  \#\mathrm{Fix}(g_6^3)=8 .
\]
The quotient $S/\langle g_6\rangle$ has two stabiliser-$\Z/6$ orbits,
two stabiliser-$\Z/3$ orbits, and two stabiliser-$\Z/2$ orbits. Its
Du-Val type is
\[
  2A_5+2A_2+2A_1 .
\]
The topological quotient has
\[
  e(S/\Z_6)=\frac{24+2+6+8+6+2}{6}=8.
\]
The exceptional curves in the crepant resolution contribute
$2\cdot5+2\cdot2+2\cdot1=16$, so the resolved surface again has Euler
characteristic $24$.

Let $r_6=\mathrm{rk}\,H^2(S,\mathbb{Z})^{g_6}$. The traces
\[
  \mathrm{tr}(g_6\mid H^2(S))=0,\qquad
  \mathrm{tr}(g_6^2\mid H^2(S))=4,\qquad
  \mathrm{tr}(g_6^3\mid H^2(S))=6
\]
give $r_6=6$ by decomposing $H^2(S,\mathbb{C})$ into order-six
eigenspaces. Therefore the invariant Mukai-lattice rank is
\[
  \mathrm{rk}\,\widetilde H(S,\mathbb{Z})^{g_6}=8 .
\]

The umbral Niemeier anchor selected by the $6A$ CHL character is
$D_4^{\oplus6}$, root system $6D_4$, with umbral group
$3.\mathrm{Sym}_6$ of order $2160$. The Niemeier lattice $A_5^4D_4$ has
outer group of order $144$ and belongs to a different Coxeter-six row.

\paragraph{Conditional denominator-to-DT assertion at level six.}
Assume a level-six Lorentzian BKM datum with root lattice, Weyl chamber,
real roots, imaginary roots, parity, multiplicities, and denominator
identity is supplied, with primitive automorphic product $\FCHL_6$ and
with root multiplicities
\[
\mult_{\mathfrak{g}^{(6)}}(n\delta+\ell z+mw)
\;=\;c^{(g_6)}_{K3}\!\bigl(4nm-\ell^2,\ell\bigr)
\]
in the positive chamber. Assume also a reduced orbifold-DT gluing theorem
identifying the local $A_5$, $A_2$, and $A_1$ vertex factors, the reduced
$S\times E$ cosection, and the primitive CHL denominator normalisation.
Under these hypotheses,
\begin{equation}
\DTred\bigl(\mathcal{X}^{(6)}\bigr)(\tau,\sigma,z)
\;=\;-\,\frac{C_6(\tau,z)}{\bigl(\FCHL_{6}(\tau,\sigma,z)\bigr)^2},
\qquad
C_6=2\phi^{(g_6)}_{0,1},\quad
\mathrm{wt}\,\FCHL_6=c_6(0)/2=1 .
\label{wn:eq:N6-DTred}
\end{equation}
Here $\phi^{(g_6)}_{0,1}\in J_{0,1}(\Gamma_0(6))$ is the $6A$-twined K3
weak Jacobi form and $c_6(0)=2$.

\paragraph{Enumerative input.}
Bryan--Steinberg 2014 proves the equivariant DT crepant-resolution
correspondence for local CY$_3$ orbifolds under finite abelian
symplectic actions preserving the holomorphic three-form.
Bryan--Cadman--Young 2013 computes the local cyclic $A$-type orbifold
topological vertex. Jun Li's degeneration formula supplies the
fibrewise gluing formalism. Oberdieck's reduced theory for $S\times E$
controls the level-one reduced $K3\times E$ scalar. The conditional
assertions above add the global compatibility between these local
ingredients and the primitive CHL products $\FCHL_4,\FCHL_6$.

References: Atiyah--Bott 1968 \emph{Ann.\ of Math.}; Nikulin 1979
\emph{Trudy Moscow Math.\ Soc.}~$38$; Mukai 1988
\emph{Invent.\ Math.}~$94$; Hashimoto 2012
\emph{Tohoku Math.\ J.}~$64$; Bryan--Steinberg 2014
\emph{Duke Math.\ J.}~$168$; Bryan--Cadman--Young 2013
\emph{J.\ Amer.\ Math.\ Soc.}~$28$; Jun Li 2001
\emph{J.\ Differential Geom.}~$60$; Oberdieck 2018
\emph{Duke Math.\ J.}~$167$; Gritsenko--Nikulin 1995 Part~II;
Borcherds 1998 \emph{Invent.\ Math.}~$132$; Gritsenko 1999
\emph{St.\ Petersburg Math.\ J.}~$11$; Cl\'ery--Gritsenko 2008;
Cheng--Duncan--Harvey 2014 \texttt{arXiv:1307.5793}.
